\section{The displayed extremizer}\label{sec:witness}
\subsection{Normalization, Gram matrix, and orthogonal rows}
Let $H_d=(\omega^{e_{ij}})$ be the dephased matrix in \cref{app:matrix}; its first row and first column are all ones. Dephasing changes the determinant by a unit, not its absolute value. Exact Eisenstein Bareiss elimination gives
\begin{equation}\label{eq:dephased-det}
 \det H_d=-362797056-604661760\omega,
 \qquad \N(\det H_d)=B.
\end{equation}
The same value is obtained by computing the determinant of its Gram matrix.

After the explicit row switching and ordering in \cref{app:matrix}, the Gram becomes
\begin{equation}\label{eq:canonical-gram}
 G_c=
 \begin{pmatrix}
 I_7\otimes C&\delta\one_{14}\\
 \bar\delta\one_{14}^T&15
 \end{pmatrix},\qquad
 C=\begin{pmatrix}15&3\\3&15\end{pmatrix}.
\end{equation}
Thus its support consists of seven triangles sharing one vertex. There are seven off-diagonal pairs of norm 9, fourteen of norm 3, and 84 zero pairs; hence $Q=105$. Its spectrum is
\begin{equation}\label{eq:witness-spectrum}
 12\ (7\text{ times}),\qquad18\ (6\text{ times}),\qquad
 \frac{33+\sqrt{177}}2,\quad\frac{33-\sqrt{177}}2.
\end{equation}
Indeed, pair differences give the eigenvalue 12; the six-dimensional zero-sum space of pair sums gives 18; and the remaining two-dimensional block has trace 33 and determinant 228. Consequently,
\[
 \chi_{G_c}(x)=(x-12)^7(x-18)^6(x^2-33x+228),\qquad
 \det G_c=12^7\,18^6\,228=B.
\]
This is a short algebraic verification of the benchmark determinant, independent of the elimination calculation.

The largest mutually orthogonal subset of rows of this particular matrix has size seven. It consists of one endpoint from each matching edge, never the common vertex. There are exactly $2^7=128$ such subsets. Each noncentral row is orthogonal to 12 other rows, whereas the central row is orthogonal to none. These are different statistics from the global orthogonality capacity $M_3(15)=9$. The column Gram has the same support and norm statistics, and the largest orthogonal column subsets likewise have size seven. The dephased row and column Grams are not equal; dephasing does not imply normality.

\subsection{Automorphisms}
We use the following fixed convention:
\[
 \Aut_{\muT}(H)=\{(M,N):M,N\text{ are monomial with nonzero entries in }\muT,
                         \ MHN^*=H\}.
\]
It contains the central subgroup $\{(\omega^jI,\omega^jI):j=0,1,2\}$, acting trivially on entries. The quotient by this subgroup is called the projective monomial automorphism group here. Complex conjugation and transposition are not included.

\begin{proposition}\label{prop:automorphisms}
For the displayed extremizer,
\[
 |\Aut_{\muT}(H_d)|=1008,
 \qquad \Aut_{\muT}(H_d)\cong C_6\times\GL(3,2).
\]
Its projective monomial automorphism group has order 336 and is isomorphic to $C_2\times\GL(3,2)$.
\end{proposition}
\begin{proof}
Use the equivalent representative $H_c$ whose Gram is \eqref{eq:canonical-gram} and whose last row is all ones; its exponent matrix and the switching are included in the data. Every row-monomial Gram stabilizer fixes the unique degree-14 vertex and permutes the seven pairs, possibly interchanging their endpoints. Since all hub entries equal $\delta$, its row phases must be equal. Modulo the central scalar subgroup, the row action is therefore one of exactly $7!2^7=645120$ permutations.

For such a permutation, the last row of $H_c$ forces all compensating column phases to be one. A compensating column permutation exists precisely when the permuted columns equal the original column set. The columns are distinct because $H_c$ is nonsingular, so it is unique. Exhaustive integer comparison of these 645120 row permutations gives exactly 336 pairs. The program verifies all 225 entries for every retained pair.

The resulting permutation group has center of order two, generated by interchanging the endpoints of all seven pairs, and derived subgroup of order 168. Their intersection is trivial. On the seven pairs, the derived subgroup acts faithfully and preserves the seven triples
\[
 123,\quad145,\quad167,\quad246,\quad257,\quad347,\quad356.
\]
Every pair of points lies in exactly one triple, so this is the Fano plane. The program independently checks all $7!$ permutations of the seven points and finds 168 preserving these triples. Its automorphism group is $\GL(3,2)$: label the points by the nonzero vectors of $\mathbb F_2^3$, with lines the triples summing to zero; a line-preserving permutation is determined by the images of a basis and preserves addition. Thus the projective group is $C_2\times\GL(3,2)$. The representatives just obtained are ordinary permutation pairs, so the scalar $C_3$ splits off in the full monomial group. This gives $C_3\times C_2\times\GL(3,2)$, as asserted.
\end{proof}
The much larger monomial stabilizer of $G_c$ alone must not be confused with the stabilizer of $H_c$. Most permutations of its seven pairs do not preserve the column set of a factor. Likewise, determining the automorphism group of one maximizer does not establish uniqueness of its Hadamard equivalence class. The search above discards $D=B$ when proving the upper bound and has not classified all equality cases.
